\section{Introduction}
Execution feedback can be learned inside a code policy or interpreted by a separate module. Reinforcement learning from execution results offers a direct route to improving iterative synthesis~\citep{rlef}, but it requires actor training and couples the learned behavior to that policy. Prompting with an expanding transcript offers another route~\citep{selfdebug}, while repeatedly exposing the actor to evidence that can be redundant, ambiguous, or specific to an obsolete program version. Neither observation implies that a compact module will work better. It motivates a controlled question: can a frozen actor exploit feedback through an explicit learned state whose predictions can be evaluated independently?

We propose Particle-Belief Program Filtering (PBPF), which separates diagnostic inference from repair generation. A learned belief module consumes semantic representations of the task, code, code change, test, and observed execution outcome. It maintains multiple weighted latent explanations for each candidate version and predicts future outcomes. The code actor supplies frozen features and remains frozen during the primary repair-conditioning stage; a small projector maps a sampled latent to soft-prefix tokens, drawing on continuous prompting~\citep{li2021prefix}. This is parameter-efficient actor conditioning, not training-free repair: belief estimation, calibration, and the conditioning interface all require training data and computation.

The state has a precise lifecycle. Tests executed on unchanged source code update likelihood weights without inventing a new latent transition. A newly patched version inherits only from its selected parent through a corrected proposal. Separate candidate versions do not exchange posterior mass. At generation time one component is selected and held throughout the continuation, so several plausible diagnoses need not be collapsed into a mean or interleaved token by token. The probabilistic machinery follows established filtering and resample--move ideas~\citep{delmoral2006,gilks2001}; the proposed interface connects these ideas to candidate-specific repair.

This decomposition makes stronger controls possible. Matched deterministic memories test whether the module benefits from more parameters rather than a posterior. Evidence shuffling and wrong-candidate feedback test whether its state tracks the relevant execution history. A shared bank of initial programs and explicit actor-call accounting distinguish diagnostic computation from additional program search. REx~\citep{rex} and a paper-spec rollout-particle adaptation~\citep{roulette} provide controls for allocation and inference-time search, while trace-rich debugging~\citep{ldb} belongs to a separate information regime.

The study is prospective: no accuracy or efficiency benefit has been established. Cross-model replication retrains model-specific modules under locked choices rather than assuming that soft prefixes transfer unchanged. The three proposed contributions are:
\begin{itemize}
\item A compact candidate-specific latent belief model with calibrated future-outcome predictions and explicit particle-state semantics, tested independently of the repair actor.
\item A frozen-actor conditioning procedure that learns a soft-prefix interface from whole-sequence mixtures and preserves one diagnosis throughout each repair, with all training costs disclosed.
\item An information- and generation-matched evaluation spanning prediction, repair, causal controls, and locked model replications, with separate ledgers for actor work, belief work, execution, and resources.
\end{itemize}
